\documentclass[10pt,a4paper]{amsart}
\usepackage[margin=19mm]{geometry}
\usepackage{amsmath,amssymb,mathtools}
\usepackage[scr=rsfs,cal=euler]{mathalfa}
\usepackage{xfrac}
\usepackage[unicode,hidelinks,bookmarks=false]{hyperref}
\newcommand{\ie}{i.e.\ }
\newcommand{\cf}{cf.\ }
\newcommand{\resp}{respectively\ }
\newcommand{\Subsection}{\subsection}
\providecommand{\nobreakdash}{\nobreak\hbox{-}}
\setlength{\emergencystretch}{2em}
\multlinegap0pt
\usepackage{bm}
\usepackage{bbm}
\usepackage{dsfont}
\usepackage{xspace}
\usepackage{cancel}
\usepackage{mathtools}
\usepackage{amsthm}
\makeatletter
\@ifundefined{theorem}{\newtheorem{theorem}{Theorem}}{}
\makeatother
\title[Hard labels sampled from sparse targets mislead rotation inv]{Hard labels sampled from sparse targets mislead rotation invariant algorithms}
\author{Avrajit Ghosh, Bin Yu, Manfred Warmuth, Peter Bartlett}
\date{Source: arXiv:2603.20967v2 (CC BY 4.0)}
\begin{document}
\maketitle

\noindent\emph{Source and licence.} Hard labels sampled from sparse targets mislead rotation invariant algorithms. Version arXiv:2603.20967v2 (CC BY 4.0), distributed under the Creative Commons Attribution 4.0 International licence, as recorded in the arXiv licence metadata for this exact version and shown on the abstract page https://arxiv.org/abs/2603.20967v2. Full rights notice: licenses/arxiv-2603.20967-CC-BY-4.0.txt.\par\smallskip

\noindent\textit{Editorial scope.} Counted unit: the Theorem whose source label is \texttt{pop-grad-restate} in the article (source0.tex lines 2297--2343, source label \texttt{pop-grad-restate}), delivered together with the article's own complete original proof of it (source0.tex lines 2345--2550, one \texttt{proof} environment of 206 lines). No mathematics is written, repaired or re-derived: statement and proof are the article's own text line for line. Required context retained uncounted: sparse-data (lines 2135--2138). Every definition, display and label of those lines is retained unchanged. Omitted from this package and not redistributed: 1--2134, 2139--2296, 2344--2344, 2551--4756. Nothing inside a retained excerpt is omitted or altered: every retained body is delivered exactly as the article's lines are written, so no omission receipt of any kind accompanies this package. Citations: the retained excerpts carry no citation command at all, so no bibliography entry of the article is transcribed and no reference is redistributed. Reference adaptation: none was needed and none was made; every reference inside the delivered unit resolves inside it, so no unresolved reference or citation is produced. Printed slips kept literal: no printed word, symbol, hypothesis or number of the counted statement or of its proof was altered. Layout and class: the article's own class is \texttt{article} with packages including microtype, graphicx, subcaption, booktabs, xcolor, hyperref, enumitem, icml2026, amsmath, amssymb, mathtools, amsthm, tablefootnote, tikz; the wrapper is the lane's pinned \texttt{amsart} editorial wrapper at 10pt on A4, which restates the theorem-like environments the retained text needs and adds the article's own macro definitions that the retained text needs (none), with extra wrapper packages (bm, bbm, dsfont, xspace, cancel, mathtools). The delivered numbering is the unit's own. Assets: no retained span contains a figure environment, a graphics-inclusion command, a PDF-page inclusion, a table or a TikZ picture; no image file is copied into this package, the package declares no asset dependency and no image file is redistributed with it. Licence: version arXiv:2603.20967v2 of this article is distributed under the Creative Commons Attribution 4.0 International licence, as recorded in the arXiv licence metadata for this exact version and shown on the abstract page https://arxiv.org/abs/2603.20967v2. A full rights notice is delivered at licenses/arxiv-2603.20967-CC-BY-4.0.txt. Characters: every retained excerpt is pure ASCII, so the delivered engine, XeLaTeX with its default Latin Modern text font, reports no missing character.
\begin{align}
\label{sparse-data}
    \mathbb{P}(y_i = 1 \mid \mathbf{x}_i) = \sigma(\langle \mathbf{x}_i, \mathbf{w}^* \rangle)
\end{align}

\begin{theorem}
\label{pop-grad-restate}
Define the population gradient $\mathbf g(\mathbf w):=\nabla \mathcal L(\mathbf w)$ and empirical gradient
$\widehat{\mathbf g}_n(\mathbf w):=\nabla \widehat{\mathcal L}(\mathbf w)$.
Let $S:=\mathrm{supp}(\mathbf w^\star)$ and denote by $\mathbf x_S$ the subvector of $\mathbf x$ on $S$, denoted from the sparse data-generation model~\eqref{sparse-data}. 

(i) Exact population gradient:
For every $\mathbf w\in\mathbb R^d$ and each coordinate $j\in[d]$,
\begin{align}
g_j(\mathbf w)
=
\mathbb E_{\mathbf x}\!\left[
x_j\Big(\sigma(\mathbf x^\top \mathbf w)-\sigma(\mathbf x_S^\top \mathbf w_S^\star)\Big)
\right].
\label{eq:pop_grad_exact-restate}
\end{align}
Moreover, with $\sigma'(t)=\sigma(t)\bigl(1-\sigma(t)\bigr)$ 
\begin{align}
    a(\mathbf w):=\mathbb E_{\mathbf x}\!\left[\sigma'(\mathbf x^\top \mathbf w)\right],
\quad
a^\star:=\mathbb E_{\mathbf x}\!\left[\sigma'(\mathbf x_S^\top \mathbf w_S^\star)\right]
\end{align}
 we have the gradient coordinate-wise forms by Stein's lemma:
\begin{align}
g_j(\mathbf w) &= w_j\, a(\mathbf w), \qquad && j\in S^c,
\label{eq:pop_grad_inactive-restate}
\\
g_j(\mathbf w) &= w_j\, a(\mathbf w) \;-\; w_j^\star\, a^\star, \qquad && j\in S.
\label{eq:pop_grad_active-restate}
\end{align}

\medskip
(ii) Sampling noise:
Define the gradient noise vector
\[
\boldsymbol\zeta(\mathbf w):=\widehat{\mathbf g}_n(\mathbf w)-\mathbf g(\mathbf w)\in\mathbb R^d,
\qquad
\zeta_j(\mathbf w):=\mathbf e_j^\top \boldsymbol\zeta(\mathbf w).
\]
Assume that for the fixed $\mathbf w$ of interest, each $\zeta_j(\mathbf w)$ is the average of $n$ i.i.d.\ centered sub-Gaussian random variables. Then for any $\delta\in(0,1)$, with probability at least $1-\delta$,
\begin{align}
\max_{1\le j\le d} |\zeta_j(\mathbf w)|
\;\le\;
4\sqrt{\frac{\log(2d/\delta)}{n}}.
\label{eq:uniform_coord_noise}
\end{align}
\end{theorem}

\begin{proof}
The population logistic gradient is given by
\[
\nabla L(\boldsymbol w)=\mathbb E_{\boldsymbol x,y}\!\left[-y\boldsymbol x\,\sigma(-y\boldsymbol x^\top \boldsymbol w)\right],
\]
hence for each $j\in[d]$,
\begin{equation}
\label{eq:gj_start-restate}
g_j(\boldsymbol w)=\mathbb E_{\boldsymbol x,y}\!\left[-y x_j\,\sigma(-y\boldsymbol x^\top \boldsymbol w)\right].
\end{equation}

Apply iterated expectation using the Tower property to \eqref{eq:gj_start-restate}:
\[
g_j(\boldsymbol w)
=
\mathbb E_{\boldsymbol x}\!\left[
\mathbb E_{y\mid \boldsymbol x}\!\left[-y x_j\,\sigma(-y\boldsymbol x^\top \boldsymbol w)\mid \boldsymbol x\right]
\right].
\]
Since this is conditioned on $\boldsymbol x$, we pull out $x_{j}$:
\begin{equation}
\label{eq:pull_out_Xj}
g_j(\boldsymbol w)
=
\mathbb E_{\boldsymbol x}\!\left[
x_j\,
\mathbb E_{y\mid \boldsymbol x}\!\left[-y\,\sigma(-y\boldsymbol x^\top \boldsymbol w)\mid \boldsymbol x\right]
\right].
\end{equation}

Writing $p(\boldsymbol x):=\mathbb P(y=1\mid \boldsymbol x)$ and since $y\in\{-1,+1\}$,
\begin{align}
\mathbb E_{y\mid \boldsymbol x}\!\left[-y\,\sigma(-y \boldsymbol x^\top \boldsymbol w)\mid \boldsymbol x\right]
&=
(-1)\sigma(-\boldsymbol x^\top \boldsymbol w)\,p(\boldsymbol x)
+
(+1)\sigma(\boldsymbol x^\top \boldsymbol w)\,(1-p(\boldsymbol x)) \nonumber\\
&=
-\sigma(-\boldsymbol x^\top \boldsymbol w)p(\boldsymbol x)+\sigma(\boldsymbol x^\top \boldsymbol w)\big(1-p(\boldsymbol x)\big).
\label{eq:cond_exp_raw}
\end{align}
Using $\sigma(-\boldsymbol x^\top \boldsymbol w)=1-\sigma(\boldsymbol x^\top \boldsymbol w)$,
\begin{align}
-\sigma(-\boldsymbol x^\top \boldsymbol w)p(\boldsymbol x)+\sigma(\boldsymbol x^\top \boldsymbol w)\big(1-p(\boldsymbol x)\big)
&=
-(1-\sigma(\boldsymbol x^\top \boldsymbol w))p(\boldsymbol x)+\sigma(\boldsymbol x^\top \boldsymbol w)-\sigma(z)p(\boldsymbol x) \nonumber\\
&=
\sigma(\boldsymbol x^\top \boldsymbol w)-p(\boldsymbol x).
\label{eq:cond_exp_simplified}
\end{align}
Under the logistic generative model, $p(\boldsymbol x)=\sigma(\boldsymbol x_S^\top \boldsymbol w_S^\star)$.
Therefore,
\begin{equation}
\label{eq:cond_exp_final}
\mathbb E_{y\mid \boldsymbol x}\!\left[-y\,\sigma(-y\boldsymbol x^\top \boldsymbol w)\mid \boldsymbol x\right]
=
\sigma(\boldsymbol x^\top \boldsymbol w)-\sigma(\boldsymbol x_S^\top \boldsymbol w_S^\star).
\end{equation}


Substitute \eqref{eq:cond_exp_final} into \eqref{eq:pull_out_Xj} to obtain
\[
\label{exact_pop_grad}
g_j(\boldsymbol w)
=
\mathbb E_{\boldsymbol x}\!\left[
x_j\Big(\sigma(\boldsymbol x^\top \boldsymbol w)-\sigma(\boldsymbol x_S^\top \boldsymbol w_S^\star)\Big)
\right],
\]
which proves \eqref{eq:pop_grad_exact-restate}.

Also when $j\in S^c$.
Since $x_j$ is independent of $\boldsymbol x_S$ and $\mathbb E[x_j] = 0$,
\[
\mathbb E_{\boldsymbol x}\!\left[x_j\,\sigma(\boldsymbol x_S^\top \boldsymbol w_S^\star)\right]
=
\mathbb E[x_j]\cdot \mathbb E_{\boldsymbol x_S}\!\left[\sigma(\boldsymbol x_S^\top \boldsymbol w_S^\star)\right]
=0.
\]
Thus from \eqref{eq:pop_grad_exact-restate},
\[
g_j(\boldsymbol w)=\mathbb E_{\boldsymbol x}\!\left[x_j\,\sigma(\boldsymbol x^\top \boldsymbol w)\right].
\]

\underline{\textit{Stein's Lemma reduction of exact population gradient}}

We notice that \eqref{eq:pop_grad_exact-restate} is of the form $\mathbb{E}(xf(x))$ with $x$ being a standard normal, so this can be reduced $\mathbb{E}(f'(x))$ if $f$ is smooth and differentiable. 

Furthermore in \eqref{eq:pop_grad_exact-restate}, $\mathbf{x}$ is i.i.d, hence writing $\boldsymbol x^\top \boldsymbol w = w_j x_j + Z$ where
$Z:=\sum_{k\neq j} w_k x_k$ is independent of $x_j$.
Conditioning on $Z$ and using Stein's lemma for $x_j\sim\mathcal N(0,1)$,
\[
\mathbb E\!\left[x_j\,\sigma(w_j x_j+Z)\mid Z\right]
=
w_j\,\mathbb E\!\left[\sigma'(w_j x_j+Z)\mid Z\right].
\]
So, we apply the Stein's lemma, and derive when $i\in S$:
\begin{align}
  g_j(\boldsymbol w)=  \mathbb E_{\boldsymbol x}\!\left[
x_j\Big(\sigma(\boldsymbol x^\top \boldsymbol w)-\sigma(\boldsymbol x_S^\top \boldsymbol w_S^\star)\Big)\right] = w_j\,\mathbb E_{\boldsymbol x}\!\left[\sigma'(\boldsymbol x^\top \boldsymbol w)\right] - w^{\star}_j\,\mathbb E_{\boldsymbol x}\!\left[\,\sigma'(\boldsymbol x_S^\top \boldsymbol w_S^\star)\right]
\end{align}
Since when $i \in S^{c}$, $w^{\star}_j = 0$, the gradient reduces to 
\begin{align}
  g_j(\boldsymbol w) =   \mathbb E_{\boldsymbol x}\!\left[
x_j\sigma(\boldsymbol x^\top \boldsymbol w)\right] = w_j\,\mathbb E_{\boldsymbol x}\!\left[\sigma'(\boldsymbol x^\top \boldsymbol w)\right] 
\end{align}

\underline{\textit{Coordinate-wise sampling noise}}





Assume $\boldsymbol x\sim\mathcal N(\boldsymbol 0,\boldsymbol I_d)$ and $Y\in\{-1,+1\}$.
 A single-sample coordinate
gradient contribution
\[
g_j(\boldsymbol w;\boldsymbol x,y):=-y\,x_j\,\sigma(-y\boldsymbol x^\top \boldsymbol w),
\]
and the centered version
\[
z_j(\boldsymbol w;\boldsymbol x,y):=g_j(\boldsymbol w;\boldsymbol x,y)-\mathbb E\!\left[g_j(\boldsymbol w;\boldsymbol x,y)\right].
\]
% Given i.i.d.\ samples $( \mathbf{x}_i,y_i)_{i=1}^n$, define the coordinate gradient noise
% \begin{align}
% \zeta_j(\boldsymbol w) =
% \big(\nabla \widehat{\mathcal L}(\boldsymbol w)\big)_j
% -
% \big(\nabla \mathcal L(\boldsymbol w)\big)_j =
% \frac{1}{n}\sum_{i=1}^n
% z_j(\boldsymbol w;\boldsymbol x_i,y_i).
% \end{align}
% Then, for any $\delta\in(0,1)$, with probability at least $1-\delta$,
% \[
% \max_{1\le j\le d}|\zeta_j(\boldsymbol w)|
% \le
% 4\sqrt{\frac{\log(2d/\delta)}{n}}.
% \]


Since $0<\sigma(\cdot)<1$ and $|y|=1$, we have
\[
|g_j(\boldsymbol w;\boldsymbol x,y)|
=
|y|\,|x_j|\,\sigma(-y\boldsymbol x^\top \boldsymbol w)
\le
|x_j|.
\]
Because $x_j\sim\mathcal N(0,1)$,
\[
\mathbb P\!\left(|g_j(\boldsymbol w;\boldsymbol x,Y)|\ge t\right)
\le
\mathbb P(|x_j|\ge t)
\le
2e^{-t^2/2},
\qquad \forall\,t\ge 0.
\]

The tail bound above implies that $g_j(\boldsymbol w;\boldsymbol x,y)$ is sub-Gaussian
(with an absolute-constant proxy). Concretely, one may take the mgf bound
\[
\mathbb E\!\left[e^{\lambda g_j(\boldsymbol w;\boldsymbol x,y)}\right]
\le
e^{2\lambda^2},
\qquad \forall\,\lambda\in\mathbb R,
\]
so $g_j(\boldsymbol w;\boldsymbol x,y)$ is sub-Gaussian with variance proxy $4$.

Centering preserves sub-Gaussianity up to an absolute constant. In particular, one can take
\[
\mathbb E\!\left[e^{\lambda z_j(\boldsymbol w;\boldsymbol x,y)}\right]
\le
e^{4\lambda^2}
=
\exp\!\left(\frac{8\lambda^2}{2}\right),
\qquad \forall\,\lambda\in\mathbb R,
\]
so $z_j(\boldsymbol w;\boldsymbol x,y)$ is sub-Gaussian with variance proxy $8$.

Since $(z_j(\boldsymbol w;\boldsymbol x_i,y_i))_{i=1}^n$ are i.i.d.\ centered sub-Gaussian
with variance proxy $8$, their average satisfies: for any $\delta\in(0,1)$,
\[
\mathbb P\!\left(
|z_j(\boldsymbol w)|
\ge
\sqrt{\frac{16\log(2/\delta)}{n}}
\right)
\le
\delta.
\]
Equivalently, with probability at least $1-\delta$,
\[
|z_j(\boldsymbol w)|
\le
4\sqrt{\frac{\log(2/\delta)}{n}}.
\]

Apply a union bound over $j=1,\dots,d$ (replace $\delta$ by $\delta/d$ in Step 4):
with probability at least $1-\delta$,
\[
\max_{1\le j\le d}|z_j(\boldsymbol w)|
\le
4\sqrt{\frac{\log(2d/\delta)}{n}}.
\]
We denote this event as $\mathcal{E}$. 
\end{proof}

\end{document}
